\documentclass[12pt]{article}
\usepackage{amsmath,amssymb}

\begin{document}
\section*{{2019 AIME I Problems/Problem 1}}
Source: AoPS Wiki

\subsection*{Solution}
Let's express the number in terms of $10^n$ . We can obtain $(10-1)+(10^2-1)+(10^3-1)+\cdots+(10^{321}-1)$ . By the commutative and associative property, we can group it into $(10+10^2+10^3+\cdots+10^{321})-321$ . We know the former will yield $1111....10$ , so we only have to figure out what the last few digits are. There are currently $321$ 1's. We know the last four digits are $1110$ , and that the others will not be affected if we subtract $321$ . If we do so, we get that $1110-321=789$ . This method will remove three $1$ 's, and add a $7$ , $8$ and $9$ . Therefore, the sum of the digits is $(321-3)+7+8+9=\boxed{342}$ . -eric2020
-another Eric in 2020 A similar and simpler way to consider the initial manipulations is to observe that adding $1$ to each term results in $(10+100+... 10^{320}+10^{321})$ . There are $321$ terms, so it becomes $11...0$ , where there are $322$ digits in $11...0$ . Then, subtract the $321$ you initially added. ~ BJHHar

\end{document}
